\section*{Discussion}
\label{sec:discussion}
\subsection*{Domain-specific tools}
AIRVIC detected CPE in 21 of the 22 images in the description set and in 100 of the 101 total images, resulting in a high false-positive rate of 0.909 and an overall accuracy of only 0.545. The test images are visually similar to the Vero-cell images used in AIRVIC’s original training set, as demonstrated by table \ref{tab:airvic-image-comparison}. Despite this similarity, the model failed to generalize effectively to the present dataset. Although direct access to the AIRVIC training data was not available, it is possible that the generally low confluency of the current dataset was not well represented in the original training set.

\clearpage
\begin{landscape}
\begin{table}[p]
\centering
\small
\setlength{\tabcolsep}{7pt}
\renewcommand{\arraystretch}{1.55}
\caption{Comparison of image acquisition and processing parameters between the AIRVIC training dataset \cite{airvic} and the current Vero cell culture dataset \cite{mathews2025}.}
\label{tab:airvic-image-comparison}
\vspace{10pt}

\begin{tabular}{p{4.2cm} p{5.3cm} p{5.3cm} p{9.0cm}}
\toprule
\textbf{Aspect} & \textbf{AIRVIC dataset} & \textbf{Vero cell
culture dataset} & \textbf{Potential impact on AIRVIC accuracy} \\
\midrule
\arrayrulecolor{gray!30}
Microscope system & ZEISS Axio Observer 5 widefield microscope & EVOS\textsuperscript{\textregistered} FL Auto system (Thermo Fisher, Cat. no. AMAFD1000) & Moderate — model trained exclusively on a single microscope brand and optical path; cross-platform generalization has not been validated \\
\hline
Objective magnification & 5×, 10×, 20×, 40× (phase contrast) & 10× and 20× (AMG LPlan FL PH objectives) & Low — strong overlap with the optimal training magnifications \\
\hline
Contrast / Illumination & Phase contrast only (Ph1/Ph2 phase rings) & Phase contrast only & Low — excellent match with training conditions \\
\hline
Camera & ZEISS Axiocam 506 color CMOS (2752 × 2208 px) & High-sensitivity CMOS (EVOS) & Moderate — differences in sensor response, color balance, and pixel characteristics possible \\
\hline
Original file format & JPEG (24-bit color) & TIFF & Negligible \\
\hline
Final format submitted to AIRVIC & JPEG (internally converted to grayscale) & PNG (converted from TIFF, color preserved) & Low–Moderate — model expects grayscale-like input; absence of explicit grayscale conversion may introduce minor color-related bias \\
\hline
Acquisition software & ZEISS ZEN lite 2.3 & DIBImage v1.0.0 (rev. 462) / Evos software & Low \\
\arrayrulecolor{black}
\bottomrule
\end{tabular}

\end{table}
\end{landscape}
\clearpage

The three pre-trained DVICE models performed similarly poorly, yielding a high false-positive rate of 0.818 and an overall accuracy of only 0.545. When the average probability across the three models was used, DVICE achieved a maximum accuracy of 0.5909 (AUROC = 0.4132) even after threshold optimization. The models produced many false positives on CPE-negative images while missing some true positives. It should be noted that the DVICE training set consisted of transmitted-light microscope images, whereas the current dataset consists of phase-contrast images (see Table \ref{tab:dvice-image-comparison}). This difference in imaging modality alone may account for the poor performance of the DVICE models in this study. The low confluency of the current dataset is unlikely to explain the poor performance, since the DVICE authors explicitly state that their process enabled ``the network to learn that low cell confluency is not a defining hallmark of viral infection state'' \cite{dvice}.

\clearpage
\begin{landscape}
\begin{table}[p]
\centering
\small
\setlength{\tabcolsep}{7pt}
\renewcommand{\arraystretch}{1.55}
\caption{Comparison of image acquisition and processing parameters between the DVICE training dataset\cite{dvice} and the current Vero cell culture dataset\cite{mathews2025}.}
\label{tab:dvice-image-comparison}
\vspace{10pt}

\begin{tabular}{p{4.2cm} p{5.3cm} p{5.3cm} p{9.0cm}}
\toprule
\textbf{Aspect} & \textbf{DVICE expectation (training data)} & \textbf{Current dataset (this study)} & \textbf{Potential impact on DVICE accuracy} \\
\midrule
\arrayrulecolor{gray!30}
Microscope system & ImageXpress Micro Confocal (IXM-C, Molecular Devices); Cytation 5 (validation) & EVOS\textsuperscript{\textregistered} FL Auto system (Thermo Fisher, Cat. no. AMAFD1000) & Moderate — model trained on a different high-throughput confocal system; cross-platform generalization not validated \\
\hline
Objective magnification & 4× air objective & 10× and 20× (AMG LPlan FL PH objectives) & Moderate — lower magnification and different optical setup than training data \\
\hline
Contrast / Illumination & Transmitted light (TL), label-free & Phase contrast only & High — fundamental difference in contrast mechanism (brightfield TL vs phase contrast) \\
\hline
Camera / Resolution & 2048 × 2048 px (16-bit) or 1992 × 1992 px (Cytation 5) & High-sensitivity CMOS (EVOS) & Moderate — differences in sensor response, bit depth, and pixel characteristics \\
\hline
Original file format & 16-bit raw (proprietary) & TIFF & Negligible \\
\hline
Final format / preprocessing & Rescaled to 8-bit PNG + resized to 224 × 224 px (RGB) & PNG (converted from TIFF, color preserved) & Moderate — model expects specific 224 × 224 RGB input after min-max normalization; current pipeline does not match this preprocessing \\
\hline
Cell line & VeroE6 / VeroE6-TMPRSS2 (SARS-CoV-2) + multiple human lines & Vero (standard) & Low — cell line is comparable \\
\hline
Acquisition timing & 7 days post-infection (dpi) & Variable (multiple passages) & Moderate — CPE morphology may differ due to timing and passage number \\
\arrayrulecolor{black}
\bottomrule
\end{tabular}

\end{table}
\end{landscape}
\clearpage

Cellpose similarly classified nearly all images as CPE-positive (accuracy of 0.500 at the default threshold of 0.5, with a maximum accuracy of 0.5909). Consequently, the morphological-proxy approach did not provide useful discrimination in this dataset.

\subsection*{Multimodal AI tools}
The four general-purpose multimodal models exhibited macro-averaged accuracies between 0.542 and 0.709. Gemini achieved the highest overall accuracy (0.709) but at the cost of a very high false-negative rate, effectively behaving like a conservative ``no-CPE'' classifier; this performance was still lower than that of the simple always-negative baseline model. ChatGPT, Claude, and Grok showed poor performance (accuracies 0.549--0.559) with substantial rates of both false positives and false negatives. None of the multimodal models reached a level of reliability that would support their use as objective CPE detection tools for this dataset.

\subsection*{Limitations}
The ground-truth labels are based on narrative descriptions provided by the CRO that were not generated under a CPE-specific protocol. The description set is small (\(n = 22\)), and the remaining 83 images could not be evaluated quantitatively. Domain-specific tools were tested exactly as publicly available, with no retraining or fine-tuning performed. Prompt engineering for the multimodal models, although extensive, remains inherently model-dependent.

Furthermore, the images from the experimental stage exhibited relatively low cell confluency compared with typical culture conditions under which CPE is normally observed. It is possible that the training sets of the domain-specific models did not include images with similarly low confluency.

\subsection*{Implications}
None of the evaluated AI tools demonstrated sufficient accuracy or robustness to replace or act as an objective second opinion for human expert interpretation of CPE in Vero-cell cultures under the conditions tested. As a consequence, an objective assessment of CPE presence could not be obtained for the remaining 83 images of the EXP stage. Future work may benefit from larger, prospectively labeled datasets and from domain-specific fine-tuning of multimodal models.